\documentclass[12pt]{article}
\usepackage[utf8]{inputenc}
\usepackage[T1]{fontenc}
\usepackage{hyperref}
\title{Ter\^ena Language Resources for K-12 Curriculum Development}
\author{Compiled by Claude Sonnet 5 (claude-sonnet-5) for Andr\'as Kornai}
\date{}
\begin{document}
\maketitle

\section{Introduction}
Ter\^ena (ISO 639-3 ter) is an Arawakan (Aruak) language of Mato Grosso do Sul, Brazil. Sources in the search describe it as agglutinative with VSO word order. One source says the Terena number almost 26,000. Bilingual education is guaranteed by law, but the sources describe difficulties in maintaining the language. Only excerpts were seen, and I did not open the documents.

\section{Primary Educational Resources}
No finished textbook or official list of Terena school materials was found. The dictionary below was designed for school use.
\begin{itemize}
\item A 2005 University of Bras\'ilia master's dissertation, ``O L\'exico da L\'ingua Ter\^ena: Proposta do Dicion\'ario Infantil Bilingue Ter\^ena-Portugu\^es'', proposes a children's bilingual dictionary for Terena schools. The link is to a journal article reviewing it, whose host is not identified in the result. \url{https://repositorio.ufc.br/bitstream/riufc/24108/1/2013_art_cpoliveirarvferreira.pdf}
\end{itemize}

\section{Linguistic and Reference Materials}
\begin{itemize}
\item A 2018 UFRJ master's dissertation on Terena phonology, semantics and syntax and the writing problems Terena teachers face. \url{https://pantheon.ufrj.br/handle/11422/23516}
\item A UFMS dissertation giving a phonological description of Terena (2009). \url{https://repositorio.ufms.br/jspui/bitstream/123456789/1119/1/Denise%20Silva.pdf}
\item ``L\'ingua terena'', a PDF in the Trabalho Indigenista library. I could not see its content. \url{https://biblioteca.trabalhoindigenista.org.br/wp-content/uploads/2018/06/lingua_terena.pdf}
\item An article on definiteness marking in Terena, in Cadernos do IL (UFRGS). \url{https://seer.ufrgs.br/index.php/cadernosdoil/article/download/128512/86922/547473}
\end{itemize}

\section{Religious and Community Resources}
None found.

\section{Digital and Audio Resources}
None found.

\section{Curriculum Development Context}
The search found academic work on Terena grammar, phonology and school writing, plus one proposed school dictionary. A UFMS pedagogy paper on the Bananal village in Aquidauana discusses how hard it is to maintain Terena as a mother tongue. Next steps are to read these works in full and to look for Terena teachers' materials and municipal and state indigenous education programmes in Mato Grosso do Sul.
\end{document}